\section{Discussion and Conclusion}
In this Phase 1 work, we implemented and evaluated the machine learning pipeline proposed by Morapedi and Obagbuwa \cite{morapedi2023air} for predicting $\text{PM}_{2.5}$ in South Africa. We categorized our findings under three labels: \textbf{reproduced} for the AQI classification and demo test, \textbf{approximately reproduced} for the spatial cross-validation, unmonitored estimations, and future forecasts, and \textbf{no paper numbers to compare} for the historical date split and per-city temporal evaluations. The regression models give moderate accuracy (RMSE about 26 to 29 AQI units, with a mean $\text{PM}_{2.5}$ AQI of 51.35 across cities) while exposing several methodological flaws in the published study.

\subsection{Methodological Flaws in the Base Paper}
During our reproduction, we identified four methodological issues: (1)~\textbf{Deterministic Label Leakage in AQI Classification}: Target \texttt{AQI\_Class} was derived deterministically from \texttt{Median\_PM25} using fixed cutoffs. Training models to predict a bucket derived from that same number represents label leakage. Accuracies near 100\% are trivial, as an \texttt{if-else} rule achieves identical results without training. In practice, AQI classifiers should predict future risk from past signals. (2)~\textbf{Discrepancy in Model Estimators}: The written text of Section 3.3 (v) specifies 1,000 estimators for XGBoost, whereas the printed model code in Section 3.3 (v) shows \texttt{n\_estimators=500}. We followed the written text (1,000 estimators). (3)~\textbf{Metric Mischaracterization}: In Section 5, the authors state continuous regressors ``achieved reasonably good Accuracy scores of over 0.6 and were both often correct for over 0.9 of the time''; continuous regression uses error loss, and this statement lacks mathematical definition. (4)~\textbf{Future Forecast with Identical Lag Values}: In Section 4.5 and Figure 12, sequential future rows contain identical lag values. Freezing lag columns removes autoregressive feedback. In our implementation, we forecast recursively day by day, updating lag columns dynamically.

\subsection{Limitations and Opportunities for Phase 2}
The primary limitation is reliance on univariate ground AQI readings without physical meteorology (boundary layer height, wind velocity, humidity, thermal inversions). In addition, spatial coordinate interpolation causes tree models to output identical predictions across broad regions (Kimberley, Kuruman, Upington, Vryburg). For Phase 2, we plan three improvements: (1)~\textbf{Meteorological and Satellite Integration}: ingesting ERA5 reanalysis and Sentinel-5P products; (2)~\textbf{Leakage-Free Predictive Classification}: predicting next-day risk categories from past inputs; and (3)~\textbf{Spatiotemporal Graph Architectures}: modeling regional dispersion with graph neural networks.

In conclusion, Phase 1 evaluated the findings of Morapedi and Obagbuwa, verified the conditions under which their published metrics arise, and established a clean, reproducible foundation for Phase 2.
